% =========================
% Explanation: Error Analysis / Diagnostic Plots (informal, readable)
% =========================
\subsection*{Explanation of Error Analysis / Diagnostic Plots (informal)}
This section describes what we inspect \emph{in addition to headline metrics} in order to understand why an optimizer succeeds or fails.
Diagnostic plots are essential in HPO/NAS benchmarking because two methods can end with similar final performance while exhibiting very different
search dynamics, robustness, and failure modes.

\begin{itemize}
    \item \textbf{Why diagnostics matter.}
    Summary metrics (e.g., regret at budget \(B\) or hypervolume) compress the full optimization trajectory into a single number.
    Diagnostics reveal \emph{how} performance is achieved: exploration versus exploitation behavior, stability across seeds, sensitivity to noise,
    and whether apparent improvements come from rare lucky runs or consistent progress.

    \item \textbf{Convergence curves vs.\ evaluations (single objective).}
    We plot the incumbent best-so-far curve \(f_{\min}(t)\) against evaluation count \(t\),
    aggregated across runs (median with IQR bands).
    These curves show whether a method:
    (i) improves early (high sample efficiency),
    (ii) stalls (premature convergence),
    or (iii) improves only late (benefits mainly at large budgets).
    Wide variability bands indicate instability across seeds.

    \item \textbf{Pareto-front progression (multiobjective).}
    For multiobjective problems, we visualize nondominated sets at fixed checkpoints (e.g., \(t\in\{10,25,50,B\}\))
    and track Pareto-quality metrics over time (hypervolume and/or \(\epsilon\)-indicator).
    This reveals whether improvements come from better extremes, improved coverage/diversity, or a uniform shift of the trade-off frontier.

    \item \textbf{Stability across seeds and performance distributions.}
    We plot per-run distributions (e.g., boxplots or violin plots) for primary metrics and/or show ECDFs (performance profiles).
    These diagnostics answer: \emph{how often} does a method achieve good performance?
    A method that wins only in a small fraction of runs may be less desirable than one that is consistently strong.

    \item \textbf{Compute--performance trade-offs.}
    When secondary objectives capture compute proxies (FLOPs, parameters, latency, memory, energy),
    we compare trade-off curves and recovered Pareto sets under equal budgets.
    This helps distinguish ``wins by spending more compute'' from genuinely better trade-offs.

    \item \textbf{Search behavior and diversity diagnostics (method-specific).}
    For sampling-based optimizers (e.g., P3Net), we examine indicators of search diversity and collapse:
    \begin{itemize}
        \item entropy or diversity of sampled categorical choices over time,
        \item coverage of continuous dimensions (e.g., histograms of sampled learning rates),
        \item duplicate proposal rates (how often the same \(\lambda\) is re-sampled),
        \item archive turnover (how frequently elites are replaced).
    \end{itemize}
    These plots diagnose mode collapse (over-concentration) versus unproductive exploration (diffuse sampling without improvement).

    \item \textbf{Failure mode analysis.}
    We track and summarize failure rates over time and by region of the search space:
    out-of-memory (OOM), NaNs/divergence, invalid configurations, and timeouts.
    We report whether performance differences are driven by robustness (fewer failures) or by risk-taking in unstable regions.
    We also verify that failure handling is applied consistently across methods (penalty assignment and budget counting).

    \item \textbf{Falsification outcomes (what would disconfirm the hypothesis).}
    Diagnostics support falsification: the proposed method is considered unsuccessful if we observe patterns such as:
    \begin{itemize}
        \item no improvement in convergence curves across most seeds (benefit not robust),
        \item Pareto fronts that improve only in a narrow region (poor coverage/diversity),
        \item strong mode collapse (entropy drops quickly, duplicates rise, progress stalls),
        \item gains explained mainly by increased failures in baselines (unfair stability advantage),
        \item improvements vanish under ablations (benefit not attributable to the intended mechanism).
    \end{itemize}
    These outcomes provide concrete criteria for rejecting or narrowing claims.
\end{itemize}
